\documentclass[11pt,oneside,reqno]{article}

\usepackage[a4paper,width=15cm,height=24cm]{geometry}

\usepackage{amsmath,amsthm,amssymb}
\usepackage{bbm}
\usepackage{mathrsfs}

%\usepackage[authoryear,longnamesfirst]{natbib}
%\usepackage{array}
\usepackage{natbib}
\usepackage{color}
\usepackage[dvipsnames]{xcolor}
\usepackage{tikz}
\usepackage{tikz-cd}


\usepackage{filemod}
\usepackage[breaklinks=true]{hyperref}

%\usepackage{showlabels}

% define theorem environments
\theoremstyle{plain}
\newtheorem{theorem}{Theorem}[section]
\newtheorem{proposition}[theorem]{Proposition}
\newtheorem{lemma}[theorem]{Lemma}
\newtheorem{corollary}[theorem]{Corollary}
\newtheorem{conjecture}[theorem]{Conjecture}

\theoremstyle{definition}
\newtheorem{question}[theorem]{Question}
\newtheorem{definition}[theorem]{Definition}
\newtheorem{example}[theorem]{Example}
\newtheorem{conjexample}[theorem]{Conjectural Example}
\newtheorem{remark}[theorem]{Remark}
\newtheorem{case}[theorem]{Case}
\newtheorem{condition}[theorem]{Condition}
\newtheorem{notation}[theorem]{Notation}
\newtheorem*{assumption}{Assumption}


\makeatletter

%% display environments
\def\be#1{\begin{equation*}#1\end{equation*}}
\def\ben#1{\begin{equation}#1\end{equation}}
\def\bes#1{\begin{equation*}\begin{split}#1\end{split}\end{equation*}}
\def\besn#1{\begin{equation}\begin{split}#1\end{split}\end{equation}}
\def\bg#1{\begin{gather*}#1\end{gather*}}
\def\bgn#1{\begin{gather}#1\end{gather}}
\def\bm#1{\begin{multline*}#1\end{multline*}}
\def\bmn#1{\begin{multline}#1\end{multline}}
\def\ba#1{\begin{align*}#1\end{align*}}
\def\ban#1{\begin{align}#1\end{align}}



%%% bracket commands
\def\given{\typeout{Command 'given' should only be used within bracket command}}
\newcounter{@bracketlevel}
\def\@bracketfactory#1#2#3#4#5#6{
\expandafter\def\csname#1\endcsname##1{%
\addtocounter{@bracketlevel}{1}%
\global\expandafter\let\csname @middummy\alph{@bracketlevel}\endcsname\given%
\global\def\given{\mskip#5\csname#4\endcsname\vert\mskip#6}\csname#4l\endcsname#2##1\csname#4r\endcsname#3%
\global\expandafter\let\expandafter\given\csname @middummy\alph{@bracketlevel}\endcsname
\addtocounter{@bracketlevel}{-1}}%
}
\def\bracketfactory#1#2#3{%
\@bracketfactory{#1}{#2}{#3}{relax}{1mu plus 0.25mu minus 0.25mu}{0.6mu plus 0.15mu minus 0.15mu}
\@bracketfactory{b#1}{#2}{#3}{big}{1mu plus 0.25mu minus 0.25mu}{0.6mu plus 0.15mu minus 0.15mu}
\@bracketfactory{bb#1}{#2}{#3}{Big}{2.4mu plus 0.8mu minus 0.8mu}{1.8mu plus 0.6mu minus 0.6mu}
\@bracketfactory{bbb#1}{#2}{#3}{bigg}{3.2mu plus 1mu minus 1mu}{2.4mu plus 0.75mu minus 0.75mu}
\@bracketfactory{bbbb#1}{#2}{#3}{Bigg}{4mu plus 1mu minus 1mu}{3mu plus 0.75mu minus 0.75mu}
}
\bracketfactory{clc}{\lbrace}{\rbrace}
\bracketfactory{clr}{(}{)}
\bracketfactory{cls}{[}{]}
\bracketfactory{abs}{\lvert}{\rvert}
\bracketfactory{norm}{\Vert}{\Vert}
\bracketfactory{floor}{\lfloor}{\rfloor}
\bracketfactory{ceil}{\lceil}{\rceil}
\bracketfactory{angle}{\langle}{\rangle}

%% define \sA \cA \bA \IA for all capital letters A ... Z %%%%%%%
\newcounter{ctr}\loop\stepcounter{ctr}\edef\X{\@Alph\c@ctr}%
	\expandafter\edef\csname s\X\endcsname{\noexpand\mathscr{\X}}
	\expandafter\edef\csname c\X\endcsname{\noexpand\mathcal{\X}}
	\expandafter\edef\csname b\X\endcsname{\noexpand\boldsymbol{\X}}
	\expandafter\edef\csname I\X\endcsname{\noexpand\mathbbm{\X}}
	\expandafter\edef\csname r\X\endcsname{\noexpand\mathrm{\X}}
\ifnum\thectr<26\repeat

%%%%%%%%%%%%%%%%%%%%%%%%%%%%%%%

\newcount\minute
\newcount\hour
\newcount\hourMins
\def\now{%
\minute=\time%
\hour=\time \divide \hour by 60%
\hourMins=\hour \multiply\hourMins by 60%
\advance\minute by -\hourMins%
\zeroPadTwo{\the\hour}:\zeroPadTwo{\the\minute}%
}
\def\zeroPadTwo#1{\ifnum #1<10 0\fi#1}

% Change layout
\numberwithin{equation}{section}
\allowdisplaybreaks[4]

\renewcommand\section{\@startsection {section}{1}{\z@}%
{-3.5ex \@plus -1ex \@minus -.2ex}%
{1.3ex \@plus.2ex}%
{\center\small\sc\mathversion{bold}\MakeUppercase}}

\def\subsection#1{\@startsection {subsection}{2}{0pt}%
{-3.5ex \@plus -1ex \@minus -.2ex}%
{1ex \@plus.2ex}%
{\bf\mathversion{bold}}{#1}}

\def\subsubsection#1{\@startsection{subsubsection}{3}{0pt}%
{\medskipamount}%
{-10pt}%
{\normalsize\itshape}{\kern-2.2ex. #1.}}


\def\blfootnote{\xdef\@thefnmark{}\@footnotetext}

\makeatother

\def\note#1{\par\smallskip%
\noindent%
\llap{$\boldsymbol\Longrightarrow$}%
\fbox{\vtop{\hsize=0.98\hsize\parindent=0cm\small\rm #1}}%
\rlap{$\boldsymbol\Longleftarrow$}%
\par\smallskip}

\renewcommand{\cite}{\citet}

\def\^#1{\ifmmode {\mathaccent"705E #1} \else {\accent94 #1} \fi}
\def\~#1{\ifmmode {\mathaccent"707E #1} \else {\accent"7E #1} \fi}

\edef\-#1{\noexpand\ifmmode {\noexpand\bar{#1}} \noexpand\else \-#1\noexpand\fi}
\def\>#1{\vec{#1}}
\def\.#1{\dot{#1}}
\def\wh#1{\widehat{#1}}
\def\wt#1{\widetilde{#1}}
\def\atop{\@@atop}

\renewcommand{\leq}{\leqslant}
\renewcommand{\geq}{\geqslant}
\renewcommand{\phi}{\varphi}
\newcommand{\eps}{\varepsilon}
\newcommand{\D}{\Delta}
\newcommand{\eq}{\eqref}

\newcommand{\dtv}{\mathop{d_{\mathrm{TV}}}}
\newcommand{\dw}{\mathop{d_{\mathrm{W}}}}
\newcommand{\dk}{\mathop{d_{\mathrm{K}}}}

\newcommand{\bigo}{\mathrm{O}}
\newcommand{\lito}{\mathrm{o}}
\newcommand{\hence}{\Rightarrow}
\newcommand{\toinf}{\to\infty}
\newcommand{\tozero}{\to0}
\newcommand{\longto}{\longrightarrow}
\newcommand{\Var}{\mathop{\mathrm{Var}}\nolimits}
\newcommand{\law}{\mathscr{L}}

\newcommand{\ER}{Erd\H{o}s-R\'{e}nyi}


\newcommand{\SL}{\mathrm{SL}}
\newcommand{\Exp}{\mathrm{Exp}}
\newcommand{\Ber}{\mathrm{Ber}}
\newcommand{\Geo}{\mathrm{Geo}}
\newcommand{\Beta}{\mathrm{Beta}}
\def\equ#1{^{(#1)}}
\def\zer#1{^{[#1]}}
\def\dsb#1{^{\langle#1\rangle}}
\newcommand{\eqd}{\stackrel{d}{=}}
\newcommand{\im}{\mathrm{i}}
\newcommand{\convd}{\stackrel{d}{\longrightarrow}}
\newcommand{\bu}{\mathbf{u}}
\newcommand{\bx}{\mathbf{x}}
\newcommand{\bzero}{\mathbf{0}}
\newcommand{\by}{\mathbf{y}}
\newcommand{\bw}{\mathbf{w}}
\newcommand{\te}{\mathbf{e}}
\def\sp#1{^{(#1)}}
\newcommand{\MVN}{\mathrm{MVN}}


\def\BC{\mathrm{BC}}
\DeclareMathOperator*{\argmax}{arg\,max}
\DeclareMathOperator*{\argmin}{arg\,min}

%%%%% My own commands and packages%%%%%%%%%%%%%%%%%%%

\usepackage{nomencl}
\makenomenclature

\usepackage{mwe}
\usepackage{enumitem}
\usepackage{xcolor}
\usepackage{mathtools} % for \bigtimes 


\newcommand{\R}{\mathbb{R}}
\newcommand{\N}{\mathbb{N}}
\newcommand{\E}{\mathbb{E}}
\newcommand{\Z}{\mathbb{Z}}
\newcommand{\Hn}{\mathbb{H}_{1,+}} % for nonnegative harmonics with h(x0) = 1

\renewcommand{\Var}{\operatorname{Var}}
\newcommand{\Cov}{\operatorname{Cov}}

\renewcommand{\P}{\mathbb{P}}
\newcommand{\bt}{\mathbf{t}}
\newcommand{\bs}{\mathbf{s}}
\newcommand{\Xn}{(X_n)_{n \in \mathbb{N}}}

\newcommand{\induced}{\text{ind}}

\newcommand{\red}{\textcolor{red}}
\newcommand{\pink}{\textcolor{pink}}




% indicator function
\newcommand{\ind}{\mathbf{1}}

% tilde commands
\newcommand{\tx}{\tilde{x}}
\newcommand{\tX}{\tilde{X}}
\newcommand{\ty}{\tilde{y}}
\newcommand{\tK}{\tilde{K}}
\newcommand{\tE}{\tilde{E}}



\newcommand{\zack}[1]{\textcolor{Violet}{Zack comment: #1}}
\newcommand{\daniel}[1]{\textcolor{Emerald}{Daniel comment: #1}}


%%%%%%%%%%%%%%%%%%%%%%%%%%%%%%%%%%%%%%%%%%
\begin{document}

\title{\sc\bf\large\MakeUppercase{Diffusion Limits in Entropy-Regularized Stochastic Control}}
\author{}
\date{}

\maketitle

\noindent \textbf{Goals}
\begin{itemize}
	\item Show convergence of the discrete-time entropy-regularized optimal control problem to the continuous-time exploratory formulation of \cite{WangZariphopoulouZhou2020}.
	\item Can we develop a measure-valued notion of entropy-regularized optimal control?
\end{itemize}



\newpage 
\section{Introduction}

Suppose we have some discrete-time dynamical system on Euclidean space:
%
\[X_{t+1} = X_t + b(X_t, u_t) + \sigma(X_t, u_t) \cdot \xi_t,\]
%
where $u_t \in U$ is the input of some controller and $\xi_t \overset{iid}{\sim} \mathcal{N}(0, I)$ is the noise. In classical control theory,
we have some loss function $\ell(x,u)$ and we wish to minimize the expected cummulative (discounted) loss over the entire trajectory $(X_t)_{t \in \N}$.
Namely we aim to solve:
%
\[V(x) = \min_{(u_t)_{t \in \N}} \E\left[ \sum_{t=0}^\infty \gamma^t \ell(X_t, u_t) \;\bigg\vert\; X_0 = x\right]\]
%
\noindent for some discount $\gamma \in (0,1)$. The value function $V$ and optimal control $(u_t)$ can be determined
using the standard Bellman equation.
%
\[V(x) = \inf_{u \in U} \{\ell(x,u) + \gamma \E V(x + b(x,u) + \sigma(x,u) \cdot \xi)\}.\]
%
\noindent where the expectation is taken with respect to the noise $\xi$. Defining:
%
\[Q_\gamma(x,u) := \ell(x,u) + \gamma \E V(x + b(x,u) + \sigma(x,u) \cdot \xi),\]
%
\noindent then this implies the optimal control and value function satisfy the implicit relations:
%
\begin{equation}
	u^*(x) \in \arg\min_{u \in U} Q_\gamma(x,u), \quad V(x) = Q_\gamma(x, u^*(x)).
	\label{eqn:optimal_control_unregularized}
\end{equation}
%
We can make things more interesting by allowing for the control policy to select its control $u_t$ at random,
possibly depending on the current state. Namely, suppose there is some deterministic function $\pi(\cdot \mid x) : S \to \mathcal{P}(U)$ defining a 
distribution over the set of controls for each possible state. Then we 
can define a Markov chain by:
%
\[X_{t+1} = X_t + b(X_t, U_t) + \sigma(X_t, U_t) \cdot \xi_t, \quad U_t \sim \pi(\cdot \mid X_t).\]
%
Namely, we sample a control independently at each point of time according to our policy $(\pi_t)$. This allows us to
model a few different phenomena:

\begin{enumerate}[label=(\roman*)]
    \item \textit{Imperfect Controls}: In many real-world settings, the controller does not have perfect control over the system, and the implemented actions may be subject to noise. By allowing for randomized control policies, we can model this inherent uncertainty in the control process.

    \item \textit{Adversarial or Strategic Environments}: In some settings, predictability can be exploited by an adversary or the environment. Entropy regularization encourages policies that remain sufficiently randomized, thereby mitigating worst-case or strategic responses to deterministic behavior.

    \item \textit{Exploration Versus Exploitation}: In reinforcement learning, we often want to encourage the agent to explore the state space in order to learn about the environment and the corresponding reward process. Randomized policies explicitly encode this exploration into the objective function.

    \item \textit{Robustness and Regularization}: Entropy regularization acts as a convex regularizer on the control problem, smoothing the optimization landscape and preventing degenerate (overly deterministic) solutions. This can significantly improve stability and sample efficiency in practice \cite{haarnoja2018softactorcriticoffpolicymaximum}.

    \item \textit{Information Constraints / Bounded Rationality}: Entropy regularization can also be interpreted as penalizing deviations from a reference policy via a Kullback--Leibler divergence. In this view, the controller faces a cost for using highly concentrated policies, reflecting limited information-processing capacity or a preference for staying close to a prior behavior (\cite{ortega2015informationtheoreticboundedrationality}). 
\end{enumerate}

The \textit{entropy-regularized optimal control problem} seeks to encourage control distributions
that simultaneously explore the state-space while maintaining a low cummulative loss. Its corresponding value function is defined by:
%
\[V(x) = \inf_{(\pi_t)} \E\left[ \sum_{t=0}^\infty \gamma^t \cdot (\ell(X_t, U_t) - \tau H(\pi(\cdot \mid X_t)))\;\bigg\vert\; X_0 = x\right],\]
%
where $H(\mu) := -KL(\mu \parallel \nu)$ is the Kullback-Leibler divergence with respect to some reference measure $\nu$ on $U$ and $\tau$ is a parameter controlling the regularization strength. In the case 
$U = \R^d$ and $\nu$ is the Lebesgue measure, $H(\mu)$ is the standard Shannon entropy. 
Again one-step analysis shows:
%
\begin{equation}
	V(x) = \inf_{\mu \in \mathcal{P}(U)} \left\{\int_U \left[Q_\gamma(x,u) + \tau \log\left(\frac{d\mu}{d\nu}\right)\right] \; \mu(du)  \right\}.
	\label{eqn:entropy_regularized_bellman_eqn}
\end{equation}

%
\noindent Define the partition function:
%
\[Z(x) := \int_U \exp\left(-\frac{Q_\gamma(x,u)}{\tau}\right) \; d\nu(u),\]
%
\noindent Then if $Z(x) < \infty$  the optimal policy $\pi^*$ and value function $V$ satisfy:
%
\begin{equation}
	\pi(u \mid x) \; d\nu \propto \exp\left(-\frac{Q_\gamma(x,u)}{\tau} \right), \quad V(x) = -\tau \log(Z(x)).
	\label{eqn:optimal_policy_entropy_regularized}
\end{equation}
%
\noindent As we see, the optimal policy is naturally a soft-min / Gibbs distribution with Hamiltonian $-Q_\gamma(x, \cdot ) / \tau$, and the value function
is essentially the log-partition function of this distribution. Entropy-regularization thus serves to smooths
the strict minimization problem in equation \eqref{eqn:optimal_control_unregularized}, and sending $\tau \downarrow 0$ in \eqref{eqn:optimal_policy_entropy_regularized} 
morally recovers \eqref{eqn:optimal_control_unregularized}.

 We would like to extend this notion of entropy-regularization to continuous time stochastic control. Given any fixed control sequence $(u_t)_{t \geq 0}$ we can define the diffusion:
%
\[dX_t = b(X_t, u_t) dt + \sigma(X_t, u_t) dW_t,\]
%
Under standard regularity conditions on $b, \sigma$, which we discuss more below, this SDE has a unique strong solution. Of course, sampling a control independently at each point in time is ill-defined in continuous time.
The paper \cite{WangZariphopoulouZhou2020} overcomes this obstacle by defining ``relaxed" or ``averaged" dynamics induced by a policy $\pi$. Namely, let:
%
	\[b_\pi(x) = \int_U b(x, u) \;\pi(du \mid x), \quad \sigma_\pi(x) := \sqrt{\int_U \sigma(x, u) \sigma^T(x,u) \;\pi(du \mid x)}, \]
	\[\quad \tilde{\ell}_\pi(x) := \int_U \ell(x, u) \; \pi(du \mid x).\]
%
If we let:
%
\begin{equation}
	dX^\pi_t = b_\pi(X^\pi_t) dt + \sigma_\pi(X^\pi_t) dW_t,
	\label{eqn:exploratory_dynamics}
\end{equation}
%
\noindent then we can define the value function of a given policy $\pi$ by:
%
\begin{equation}
	V_\pi(x) := \E\left[ \int_0^\infty e^{-\rho t} \cdot (\tilde{\ell}_\pi(X^\pi_t) - \tau H(\pi(\cdot \mid X^\pi_t))) dt\;\bigg\vert\; X^\pi_0 = x\right].
	\label{eqn:exploratory_value_function}
\end{equation}

\noindent where the expectation is over the Brownian motion $\{W_t\}_{t \geq 0}$. Using this, we can solve the entropy-regularized stochastic optimal control problem:
%
\begin{equation}
	V(x) := \inf_{\pi \in \cA(x)} V_\pi(x),
	\label{eqn:exploratory_optimal_value_function}
\end{equation}
%
\noindent  where $\cA(x)$ is the set of admissible controls, which possibly depends on the state $x$. \cite{WangZariphopoulouZhou2020} call this the \textit{exploratory} formulation, and motivate it as the average local dynamics 
across many realizations of the Brownian motion paths. Again, suitable regularity on $b, \sigma$ are required for the SDE \ref{eqn:exploratory_dynamics} to have a unique strong solution.


 In this short note, we show that the exploratory dynamics of \cite{WangZariphopoulouZhou2020} arise as the scaling limit of 
a natural sequence of entropy-regularized discrete-time processes that sample controls at increasingly fine intervals. This gives a possibly more natural interpretation of
entropy-regularization in continuous time as \red{todo}. Our main result is the following:

\begin{theorem}[Entropy-Regularized Diffusive Limits]
	
	\label{thm:entropy_regularized_diffusive_limits}

	Given drift $b : \R^d \times U \to \R^d$, covariance $\sigma:\R^d\times U \to \R^{d \times d}$, loss function $\ell : \R^d \times U \to \R$,
	discount $\gamma = e^{-\rho}$, and entropy-regularization parameter $\tau$, define a family of processes $(Y^h_t)_{t \geq 0}$ by:
	%
	\[X^h_{n+1} - X^h_{n} := h \cdot b(X^h_{n}, u^h_n) + \sqrt{h} \cdot \sigma(X^h_{n} , u^h_n) \cdot \xi_n,  \quad \xi_n \overset{iid}{\sim} N(0, I_d),\]
	\[Y^h_t := X^h_{\lfloor t/h \rfloor}.\]
	%
	\noindent where the control $u^h_n$ is sampled from the optimal policy $\pi^h$ in Equation \eqref{eqn:optimal_policy_entropy_regularized} for the entropy-regularized 
	control problem with parameters $(hb, \sqrt{h} \sigma,  h \ell, \gamma^h,  h\tau)$. Assume that:
	
	\begin{enumerate}
		\item [(i)] \red{aggregate proper conditions from lemmas}. 
	\end{enumerate}

	\noindent Then as $h \to 0$, we have $Y^h_t \to Y_t$ in distribution in the Skorokhod space $D([0,T], \R^d)$,
	 where $Y_t$ is \red{the optimally controlled exploratory dyamics} a diffusion process satisfying:
	%
		\begin{equation}
			dY_t = b(Y_t, \pi) dt + \sigma(Y_t, \pi) dW_t,
			\label{eqn:limiting_entropy_diffusion}
		\end{equation}
		%
	\end{theorem}

%%%%%%%%%%%%%%%%%%%%%%%%%%%%%%%%%%%%%%%%%%%%%%%%
% Allowable Controls
%%%%%%%%%%%%%%%%%%%%%%%%%%%%%%%%%%%%%%%%%%%%%%%%

\section{Allowable Controls}

In this section we discuss the conditions on the control policies $\pi$, drift $b$, and covariance $\sigma$ to ensure that the 
stochastic optimal control problem \eqref{eqn:exploratory_optimal_value_function} \red{is well-posed}. As \cite{WangZariphopoulouZhou2020} mention, any $\pi \in \cA(x)$ in the set of admissible 
controls must at least satisfy:
%
\begin{enumerate}[label=(\roman*)]
    \item For each $t \ge 0$, $\pi_t \in \mathcal{P}(U)$ almost surely.

    \item For each $A \in \mathrm{Borel}(U)$, the process
    \[
    \left\{ \int_A \pi_t(u)\,du \;:\; t \ge 0 \right\}
    \]
    is $\{\mathcal{F}_t\}$-progressively measurable.

    \item The exploratory SDE \eqref{eqn:exploratory_dynamics} admits a unique strong solution
    $X^\pi = \{X^\pi_t : t \ge 0\}$ when the control $\pi$ is applied.

    \item The value function $V^\pi(x)$ defined in \eqref{eqn:exploratory_value_function} is finite for all $x \in \R^d$.
\end{enumerate}
%

\zack{The below is more a note to myself.}\\

\noindent Condition (ii) ensures $\pi$ is measurable as a measure-valued stochastic process, and that the coefficients of \eqref{eqn:exploratory_value_function}
are progressively-measureable and hence the underlying Ito integrals are well-defined. Condition (iii) ensures the stochastic process $\tilde{X}^\pi_t$ and hence
its corresponding value function $V^\pi$ is well-defined.
 
Conditions (i)–(iv) ensure that each admissible relaxed control induces a well-defined controlled diffusion 
with finite objective value. However, these assumptions alone do not
 guarantee existence of an optimal control. For existence, one typically also requires compactness or tightness 
 of the admissible control class together with continuity and lower semicontinuity properties of the coefficients
and cost functional, so that minimizing sequences admit convergent subsequences whose limits remain admissible.

\zack{Possibly some sort of compactness is required for the space of admissible controls.}

\noindent In the case of the standard linear-quadratic controller where $U = \R^d$ and:
%
\[b(x,u) = Ax + Bu, \quad \sigma(x,u) = Cx + Du\]
%
\noindent \red{and quadratic loss}. Given policy $\pi$, let:
%
\[\mu_t := \int_U u \pi_t(du), \qquad \sigma_t^2 := \int_U (u - \mu_t)^2 \pi_t(du). \]
%
\noindent As in \cite{WangZariphopoulouZhou2020}, we will assume the following::

\begin{enumerate}[label=(\alph*)]
	\item Conditions (i) and (ii) above. 
	\item For each $t \geq 0$, $\E_\pi \left[\int_0^t (\mu_s^2 + \sigma_s^2) ds\right] < \infty$.
	\item $\lim \inf_{T \to 0} e^{-\rho T} \E (X_T^\pi)^2 = 0$.
	\item $\E \left[\int_0^\infty e^{-\rho t} |L(X^\pi_t, \pi_t)| \; dt\right] < \infty.$
\end{enumerate}

\noindent 

\zack{I guess in this paper, we are only interested in controls where $\pi_t$ is a function of the current state alone, 
while in general the results are stated for any $\cF_t$-progressively-measurable, measure-valued process $(\pi_t)$}

%%%%%%%%%%%%%%%%%%%%%%%%%%%%%%%%%%%%%%%%%%%%%%%%
% Proof of Main Result
%%%%%%%%%%%%%%%%%%%%%%%%%%%%%%%%%%%%%%%%%%%%%%%%
\section{Proof of Theorem \ref{thm:entropy_regularized_diffusive_limits}}

Recall given policy $\pi:\R^d \to \mathcal{P}(U)$, drift $b:\R^d \times U \to \R^d$ and covariance
 $\sigma:\R^d\times U \to \R^{d \times d}$, we define for each $h > 0$:
	%
	\[X^h_{n+1} - X^h_{n} := h \cdot b(X^h_{n}, u^h_n) + \sqrt{h} \cdot \sigma(X^h_{n} , u^h_n) \cdot \xi_n,  \quad Y^h_t := X^h_{\lfloor t/h \rfloor},\]
	\[\xi_n \overset{iid}{\sim} N(0, I_d), \quad u^h_n \mid \{X^h_n = x\} \sim \pi(\cdot \mid x).\]
	%
 where the controls $u_n^h$ are sampled independently of both $\xi_n$ and $X_n^h$. Heuristically, this gives the natural
 discrete-time approximation to the exploratory dynamics \eqref{eqn:exploratory_dynamics} when we sample controls independently every $h$ seconds.

We break the proof of Theorem \ref{thm:entropy_regularized_diffusive_limits} into some key lemmas. 


%%%%%%%%%%%%%%%%%%%%%%%%%%%%%%
% Fixed Policy convergence of paths
%%%%%%%%%%%%%%%%%%%%%%%%%%%%%%
\subsection{Fixed Policy Convergence of Paths}


First, we show for a fixed policy $\pi$, we get
convergence to the exploratory dynamics \eqref{eqn:exploratory_dynamics}. 

\begin{lemma}[Fixed-Policy Convergence to Exploratory Dynamics]
	Suppose $\pi$ satisfies:
	%
	\[\sup_{||y|| \leq R} \bigg| \int_U b_i(y, u) b_j(y,u) \; \pi(du \mid y) \bigg| < \infty, \quad \sup_{||y|| \leq R} \bigg|\int_U \mathrm{Tr}(\sigma\sigma^T)^{1+\delta/2} \; \pi(du \mid y) \bigg| < \infty.\]

	\noindent for some $\delta > 0$. Moreover, suppose $b_\pi(y), \sigma_\pi(y)$ are continuous 
	 and there is a unique strong solution $X^\pi_t$ to the SDE 
	\eqref{eqn:exploratory_dynamics}. If $Y^h_0 = x_h \to x$, then $Y^h_t \Rightarrow X^\pi_t$ with $X^\pi_0 = x$ as 
	$h \downarrow 0$ weakly in the Skorokhod space $D([0,T], \R^d)$ for any $T < \infty$.

	\label{lem:convergence_fixed_policy_to_exploratory_dynamics}
\end{lemma}

\begin{proof}[Proof of Lemma \ref{lem:convergence_fixed_policy_to_exploratory_dynamics}]
	
	 Let $\Delta Y^h_t := Y^h_{t+h} - Y^h_t$. Note by the integrability constraints:
	%
	\begin{align*}
		\E [\Delta Y^h_{nh} \mid Y^h_{nh} = y]  & = \E_\xi \int \left(hb(y, u) + \sqrt{h} \sigma(y,u) \xi \right) \; \pi(du \mid y)\\
		& = h \int b(y,u) \pi(du \mid y) + \sqrt{h} \int \sigma(y,u) \pi(du \mid y) \cdot \E_\xi[\xi]\\
		%
		& = h b_\pi(y)
	\end{align*}
	%
	\noindent Hence the drift always takes the correct form. Similarly we see:
	\begin{align*}
			& \E [(\Delta Y^h_{nh}) (\Delta Y^h_{nh})^T  \mid Y^h_{nh} = y] \\
			& = \E_\xi \int \left(hb(y,u) + \sqrt{h} \sigma(y,u) \xi \right) \left(hb(y,u) + \sqrt{h} \sigma(y,u) \xi \right)^T \; \pi(du \mid y) \\
				& = h^2 \int b(y,u) b(y,u)^T \pi(du \mid y) + h \int \sigma(y,u) \sigma(y,u)^T \pi(du \mid y) \\
				& \hspace{2.08in} + h^{3/2} \int b(y,u) (\sigma(y,u) \cdot \E_\xi[\xi])^T \pi(du \mid y)\\
			& =  h \sigma_\pi^2(y) + h^2 \int_U b(y,u) b(y,u)^T \pi(du \mid y).
	\end{align*}
	%
	\noindent Hence the infinitesmal variance is correct up to an $o(h)$ correction. Our assumption on $\pi$ imply this correction is bounded uniformly for $y$ in compact sets.
	Moreover, the probability of a jump larger than $\eps$ is at most: 
	%
	\begin{equation}
		\P\left(||\Delta Y^h_{nh}||_2 > \eps \mid Y^h_{nh} = y\right)  \leq \P(||b(y, U_y)|| > \epsilon / h) + \P(||\sigma(y, U_y) \xi|| > \epsilon / \sqrt{h})
		\label{eqn:no_large_jumps}
	\end{equation}

	\noindent for $U_y \sim \pi(\cdot \mid y)$. By Markov's inequality, the first term satisfies:
	%
		\[\P(||b(y, U_y)|| > \epsilon / h) \leq \frac{h^2}{\epsilon^2} \int_U ||b(y,u)||^2 \pi(du \mid y).\]
	%	
	\noindent which is $o(h)$ uniformly for $|y| \leq R$ by assumption. Defining $V_y := ||\sigma(y, U_y) \xi||$:
	%
	\begin{align*}
		\E V_y^{2+\delta} & = \int_U \E_\xi ||\sigma(y,u) \xi ||^{2+\delta} \pi(du \mid y)\\
		%
		 & \leq \E_\xi ||\xi||^{2+\delta} \cdot \int_U  ||\sigma(y,u)||_F^{2+\delta} \;  \pi(du \mid y)\\
		 %
		 & = C_\delta \int_U \mathrm{Trace}(\sigma \sigma^T)^{1+\delta/2} \;  \pi(du \mid y). 
	\end{align*}
	%
	\noindent Hence again by Markov's inequality we see:
	%
	\[ \P(||\sigma(y, U_y) \xi|| > \epsilon / \sqrt{h}) \leq \frac{h^{1+\delta/2}}{\epsilon^{2+\delta}} \cdot \E V_y^{2+\delta},\]
	%
	\noindent which is again $o(h)$ uniformly in $|y| \leq R$ by assumption. Hence by \eqref{eqn:no_large_jumps}, the probability of large 
	jumps is $o(h)$ uniformly for $|y| \leq R$. 
	
	By assumption, the limiting SDE \eqref{eqn:exploratory_dynamics} has a unique strong solution, and hence the associated
	martingale problem is well-posed. By the standard functional convergence results (e.g. Section 8.7 of \cite{durrett1996stochastic})), we have $Y^h_t \Rightarrow X^\pi_t$
	in $D([0,T], \R^d)$.
\end{proof}

\noindent By \textit{the associated martingale problem is well-posed} we mean that for the generator:

\[Lf = \langle\nabla f,  b_\pi \rangle + \frac{1}{2} \mathrm{Trace}(\sigma_\pi^T Hf \sigma_\pi) \]

\noindent of our limiting Ito process, there exists a unique law on $C([0,T], \R^d)$ so that:
%
\[M_t := f(X_t) - f(x) - \int_0^t Lf(X_s) \; ds, \quad X_0 \overset{a.s.}{=} x,\]
%
\noindent is a (local) martingale for every $f \in C^2_b(\R^d)$. Durret provides a sufficient condition for 
this in Section 8.7 of \cite{durrett1996stochastic}, condition (A):

\begin{assumption}[A]
	\label{assumption:A}
	The coefficients $\sigma_\pi^{ij}$ and $b_\pi^i$ are continuous, and for each $x$ there is a unique measure $P_x$ on $(C,\mathcal C)$ such that $P_x(X_0=x)=1$ and
	\begin{equation*}
	\begin{aligned}
	X_t^i - \int_0^t b^i(X_s)\,ds,
	\qquad
	X_t^iX_t^j - \int_0^t \sigma_\pi^{ij}(X_s)\,ds
	\end{aligned}
	\end{equation*}
	are local martingales.
\end{assumption}

The main benefit of (A) is it avoids reference
to the domain $C_b^2(\R^d)$ of test functions, and allows us to work explicitly with the limiting infinitesmal drift 
and variance. 


%%%%%%%%%%%%%%%%%%%%%%%%%%%%%%%%%%%%
% Fixed Policy convergence of Value Function
%%%%%%%%%%%%%%%%%%%%%%%%%%%%%%%%%%%%
\subsection{Fixed Policy Convergence of Value Function}
\noindent Next, we show that again for a fixed policy $\pi$, the discrete-time value functions converge uniformly on compact sets
to the analogous continuous-time value function defined by the exploratory dynamics under the proper scalings of the relevant parameters. Namely:
%
\begin{lemma}[Convergence of Value Functions]
	Given policy $\pi: \R^d \to \mathcal{P}(U)$, loss $\ell: \R^d \times U \to \R$, discount $\gamma = e^{-\rho}$, and regularization parameter $\tau$, define for each $h > 0$:
	%
	\[V^h_\pi(x) = \E\left[ \sum_{t=0}^\infty e^{-\rho h t} \cdot h g_\pi(X^h_t)\;\bigg\vert\; X^h_0 = x\right].\]
	%
	\noindent where $g_\pi(x) := \int_U \ell(x, u) + \tau \log(\pi(u \mid x)) \; \pi(u \mid x) d\nu(u)$ is the instantaneous cost under policy $\pi$ at state $x$.
	%
	\noindent Suppose: \zack{I guess we are implicitly assuming the space $U$ of controls has some metric structure.}

	\begin{enumerate}[label=(\roman*)]
		\item \textit{Lipschitz Coefficients:} The coefficients $b(x, u), \sigma(x,u)$ are Lipschitz in $x,u$:
		       \[|b(x,u) - b(y,v)| + |\sigma(x,u) - \sigma(y,v)|  \leq C |x-y| + D |u-v|.\]

			   \noindent where $|\sigma|^2 = \sum_{i,j} |\sigma_{ij}|^2$ is the Frobenius norm.
		\item \textit{Lipschitz Cost:} The instantaneous cost $g_\pi(x)$ is Lipschitz. 
		\item \textit{Lipschitz Policy:} The policy $\pi(\cdot \mid x)$ is Lipschitz in the Wasserstein space $W_2$:
				\[d_{W_2}(\pi(\cdot \mid x), \pi(\cdot \mid y)) \leq F |x-y|, \quad and \quad \int_U |u|^2 \pi(u \mid x) < \infty \quad \forall x.\]
		\item \textit{Non-vanishing Variance}: There exists $\lambda > 0$ so $\sigma_\pi^2(x) \succeq \lambda I$ for all $x$. 


		\item \textit{Sufficient Discounting:} The discount $\rho$ is large enough. Namely, there exists a finite discount $\rho$ so that the below conclusion holds.

		
	\end{enumerate}
	
	\noindent Then as $h \to 0$, we have $V^h_\pi \to V_\pi$ uniformly on compact sets, where $V_\pi$ is given by the exploratory dynamics \eqref{eqn:exploratory_value_function}. 
	\label{lem:convergence_of_value_functions}
\end{lemma}

\zack{Possibly also have to assume the moment bounds from the convergence result? As we use below the fact that we have weak convergence of the discrete chain to the limiting diffusion.}

First, observe that the above assumptions imply the relaxed coefficients $b_\pi, \sigma_\pi$ and the policy $\pi$ satisfy a natural
global growth condition. For $\pi$, note by its $W_2$-Lipschitzness:
%
\[\int_U |u|^2 \; \pi(u \mid x) \leq C(1 + |x|^2).\]	
%
\noindent Using this, note:
%
\[|b_\pi(x)| \leq \E |b(x, U_x)| \lesssim 1 + |x| + \E |U_x| \lesssim 1 + |x|,\]
\[|\sigma_\pi(x)|^2 = \int_U \mathrm{Tr}(\sigma \sigma^T) \pi(u \mid x) = \E | \sigma(x, U_x)|^2 \lesssim 1 + |x|^2.\]
%
\noindent The latter of course implies $|\sigma_\pi(x)| \lesssim 1 + |x|$. Moreover, $b_\pi(x), \sigma_\pi(x)$ are Lipschitz and locally Lipschitz respectively. For $b_\pi$, simply observe that:
 %
 \begin{align*}
	\Big|b_\pi(x) - b_\pi(y)\Big| & = \bigg|\int_U b(x,u) \pi(u \mid x) - \int_U b(y,v) \pi(v \mid y) \bigg|\\
	%
	& \leq \int_U |b(x,u) - b(y,u)| \pi(u \mid x) + \bigg|\int_U b(y,u) \pi(u \mid x) - \int_U b(y,v) \pi(v \mid y) \bigg|\\
	%
	& \leq C |x-y| + \E |b(y, U_x) - b(y, U_y)|. 
\end{align*}
 %
 \noindent for any coupling $(U_x, U_y)$ of $\pi(\cdot \mid x), \pi(\cdot \mid y)$. Using the Lipschitzness of
  $b(y, \cdot)$ and $\pi$ proves the claim. Similarly for $x,y \in K$ compact:
	\begin{align*}
		\Big| \sigma_\pi(x) - \sigma_\pi(y) \Big|^2 & \leq \bigg|\sqrt{\int_U \sigma(x,u)\sigma(x,u)^T \pi(u \mid x)} - \sqrt{\int_U \sigma(y,u)\sigma(y,u)^T \pi(u \mid y)} \bigg|^2
	\end{align*}

	\noindent which is only $\frac{1}{2}$-H\"older continuous without further assumptions as the square root map $A \mapsto A^{1/2}$ is not Lipschitz when $A$ is close to singular. Assumption (v) essentially 
	forces our relaxed covariance matrices $\sigma_\pi^2$ to be sufficiently non-singular, so they are restricted to a region where the square-root map is actually Lipschitz. Under this assumption, Cauchy-Schwarz and the Lipschitzness of $\sigma, \pi$ imply:
	%
	\begin{align*}
		& \lesssim \frac{1}{\lambda} \bigg| \int_U \sigma(x,u)\sigma(x,u)^T \pi(u \mid x) - \int_U \sigma(y,v)\sigma(y,v)^T \pi(v \mid y) \bigg|^2 \\
		%
		& \leq \lambda^{-1}\sup_{x \in K} \E |\sigma(x, U^x)|^2 \cdot \E |\sigma(x, U^x) - \sigma(y, U^y)|^2\\
		%
		& \lesssim |x-y|^2 \lambda^{-1}\sup_{x \in K} \E |\sigma(x, U^x)|^2
	\end{align*}
	%
	\noindent Again, Lipschitzness implies this supremum is finite, so $\sigma_\pi$ is locally Lipschitz.
	
	By standard SDE theory, the above ensures the SDE \eqref{eqn:exploratory_dynamics} 
	has a unique strong solution, and hence $V_\pi$ is well-defined. The above parameter choices --- $e^{-\rho h}$ for the discount, $h\ell$ for the loss, and 
	$h \tau$ for the entropy-regularization --- are essentially chosen so that $h$ steps of the discrete-time process $X^{h}_n$ corresponds to one unit of time
	in the continuous-time process.

 % Proof
\begin{proof}[Proof of Lemma \ref{lem:convergence_of_value_functions}]
	We can express $V_\pi^h$ in terms of $Y^h_t$ and its transition semigroup $P^h_t$:
	%
	\[V_\pi^h(x) = \E\left[ \sum_{t=0}^\infty e^{-\rho h t} \cdot h g_\pi(Y^h_{ht}) \;\bigg\vert\; Y^h_0 = x\right] = \int_0^\infty e^{-\rho \lfloor t/h \rfloor h} P_t^hg_\pi(x) dt.\]
	%
	\noindent Likewise if $P_t$ is the transition semigroup associated to $X_t^\pi$:
	%
	\[V_\pi(x) = \int_0^\infty e^{-\rho t} P_tg_\pi(x) dt.\]
	%
	\noindent Pointwise convergence $V^h_\pi \to V_\pi$ follows almost immediately from the weak
	convergence of $Y^h_t \Rightarrow X^\pi_t$ that we develop in Lemma \ref{lem:convergence_fixed_policy_to_exploratory_dynamics}. Simply apply this result to the functional:
	%
	\[F:D([0,T], \R^d) \to \R, \quad F(\omega) = \int_0^T e^{-\rho t} g_\pi(\omega(t)) dt.\]
	%
	The above assumptions give us sufficient control on the tails $\int_T^\infty e^{-\rho t} |P_t g_\pi(x)| dt$ then to ensure the value 
	functions converge. The main obstacle is Lemma \ref{lem:convergence_fixed_policy_to_exploratory_dynamics} gives us no control on how this convergence
	behaves as we vary the initial condition, which we will require to prove uniform convergence. The key ingredient will be a stability result
	with respect to the initial condition, which shows that trajectories diverge at most exponentially fast:
	%
	\[\E |X_t^{\pi, x} - X_t^{\pi, z}| \lesssim |x-z| e^{Ct}.\]
	%
	This stability result follows from the Lipschitzness via classical arguments, which we will reiterate below.
	We will reduce the proof of the result to a few claims. Below, for $a,b \geq 0$ let $a \lesssim b$ denote that $a \leq Cy$ for some $C \geq 0$ depending
	only on the compact set $K$, time horizon $T$, and all the relevant Lipschitz constants. In some cases, such as the tail estimates of Claim 3 below,
	the dependence on $T$ will be important, in which case we will specifically emphasize this. \\
	
	% Claim 1 %%%%%%%%%%%%%%%%%%%%%%%%%%%%%%%%%
	\noindent \textbf{Claim 1:} For every compact $K\subset \R^d$ and $T<\infty$,
    \[
		\sup_{t\in[0,T]}\sup_{x\in K}|P_t^h g_\pi(x)-P_t g_\pi(x)|\to 0
		\quad \text{as } \quad h\downarrow 0.
		\]

	\begin{proof}[Proof of Claim 1]
		Let $\{z_i\}$ be an $\eps$-net of $K$, $x \in K$, and $t$ fixed.  Note:
		%
		\[|P_t^h g_\pi(x) - P_t g_\pi(x)| \leq |P_t^h g_\pi(x) - P_t^h g_\pi(z_i)| + |P_t^h g_\pi(z_i) - P_t g_\pi(z_i)| + |P_t g_\pi(x) - P_t g_\pi(z_i)|\]
		%
		\noindent The middle and last terms are easy to handle. For the middle, note since $(Y^h_t) \Rightarrow (X^\pi_t)$ weakly in $D([0,T], \R^d)$,
		Skorokhod's representation theorem implies we can construct versions $\tilde{Y}^h \to \tilde{X}^\pi$ converging almost surely in $D([0,T], \R^d)$. Since $\tilde{X}^\pi \in C([0,T], \R^d)$ almost
		surely, this convergence thus must be almost surely uniform. Thus as $g_\pi$ is Lipschitz:
		%
		\[\sup_{t \in [0,T]}|P_t^h g_\pi(z_i) - P_t g_\pi(z_i)| \lesssim \sup_{t\in [0,T]} \E |\tilde{Y}^h_t - \tilde{X}^\pi_t| \leq   \E \left[ \sup_{t\in [0,T]}|\tilde{Y}^h_t - \tilde{X}^\pi_t| \right].\]
		%
		\noindent If we can show $Z^h := \sup_{t \in [0,T]} |\tilde{Y}^h_t - \tilde{X}^\pi_t|$ are uniformly integrable when $h \leq 1$, this tends to zero by Vitali's convergence theorem.
		 To do so, we show they are uniformly bounded in $L^2$, which will follow from Lipschitzness and our growth conditions. Note:
		\[\E Z_h^2 \leq 2\E \sup_{t \in [0,T]} |Y^h_t|^2 + 2\E \sup_{t \in [0,T]} |X_t^\pi|^2.\]
		%
		\noindent First, note Cauchy-Schwarz and the BDG inequality imply:
		%
		\begin{align*}
			\E \sup_{t \in [0,T]} |X_t^\pi|^2 & \lesssim |x_0|^2 + \E \sup_{t \in [0,T]} \left(\int_0^t b \; ds\right)^2 + \E \sup_{t \in [0,T]} \left(\int_0^t \sigma dB_s\right)^2\\
			%
			& \lesssim |x_0|^2 + T \E \int_0^T b_\pi^2 \; ds + \E \int_0^T \sigma_\pi^2 \; ds.
		\end{align*}
		%
		\noindent Since $b_\pi, \sigma_\pi$ have linear growth:
		%
		\begin{align*}
			\E \sup_{t \in [0,T]} |X_t^\pi|^2 & \lesssim |x_0|^2 + (T + 1) \int_0^T (1 + \E |X^\pi_s|^2)\; ds\\
			%
			& \lesssim 1 + |x_0|^2 + \int_0^T \E \sup_{t\in [0,s]} |X^\pi_t|^2 \; ds.
		\end{align*}
		%
		\noindent Gr\"onwall's inequality then implies $E \sup_{t \in [0,T]} |X^\pi_t|^2 \lesssim 1 + |x_0|^2$ . Similarly, for $Y^h_t$:
		%
		\begin{align*}
			\E \sup_{t \in [0,T]} |Y^h_t|^2 & = \E\max_{n \leq \lfloor T/ h\rfloor} |X_n^h|^2 \\ 
			%
			& \lesssim |x_0|^2 + h^2 \cdot \E \sup_{t \in [0,T]}   \left(\sum_{i=1}^{\lfloor T/h \rfloor} b\right)^2 + h \E \sup_{t \in [0,T]} \left(\sum_{i=1}^{\lfloor T/h \rfloor} \sigma \xi_i\right)^2,\\
			%
			& \lesssim |x_0|^2 + h T \sum_{i=1}^{\lfloor T/h \rfloor} \E b^2 + h\sum_{i=1}^{\lfloor T/h \rfloor} \E \sigma^2 \\
			%
			& \lesssim 1 + |x_0|^2 + h\sum_{i=1}^{\lfloor T/h \rfloor} \E  \max_{t \leq i} |X^h_t|^2,
		\end{align*}
		%
		\noindent by linear growth of $b, \sigma$ and $\pi$. By discrete Gr\"onwall's inequality:
		%
		\[\E\max_{n \leq \lfloor T/ h\rfloor} |X_n^h|^2 \lesssim (1+|x_0|^2)e^{Ch \lfloor T/h \rfloor} \lesssim 1 + |x_0|^2. \]
		%
		\noindent which finishes the proof of the desired uniform integrability.

		 For the last term, note our assumptions imply our limiting diffusion $X^\pi_t$ is \red{Feller continuous}, which would suffice if $g_\pi$ were bounded. 
		 Instead, we can show directly $P_tg_\pi(x)$ is Lipschitz in $x$. As $g_\pi$ is globally Lipschitz:
		%
		\[|P_t g_\pi(x) - P_t g_\pi(z)| \lesssim \sqrt{\E |X^{\pi, x}_t - X^{\pi, z}_t|^2}, \]
		%
		\noindent where $X^{h,x}_t$ denotes the unique strong solution started at initial condition $x$. Using the Lipschitzness of the relaxed coefficients, $b, \sigma$, a standard application of Ito's isometry and Gr\"onwall's inequality implies:
		%
		\[\E |X^{\pi, x}_t - X^{\pi, z}_t|^2 \lesssim |x-z|^2.\]
		%
		\noindent See for example the proof of Theorem 5.2.1 in \cite{oksendal2003sde}. \zack{The stated proof requires $\sigma$ to be globally Lipschitz rather than just locally; some truncation argument is required maybe.} This gives:
		%
		\[|P_t g_\pi(x) - P_t g_\pi(z_i)| \lesssim \epsilon.\]
		%
		Finally, we handle the first term. As $g_\pi$ is Lipschitz:
		%
		\[|P_t^h g_\pi(x) - P_t^h g_\pi(z_i)| \lesssim \E| X^{h,x}_t -  X^{h, z_i}_t|.\]
		%
		\noindent To bound this expectation, let $D_n := X^{h,x}_n -  X^{h, z}_n$ be the current difference between the coupled processes and $\Delta X^{h,x}_n := X^{h,x}_{n+1} - X^{h,z}_n$ the one-step
		increment in one of the processes. Observe for \textit{any} coupling of $(X_t^{h,x}, X_t^{h,z})$:
		%
		\begin{align*}
			\E \left[D_{n+1}^2 \mid \cF_n\right] & = D_n^2 + 2 D_n \cdot \E\left[\Delta X^{h,x}_n - \Delta X^{h,z}_n \mid \cF_n\right] + \E\left[(\Delta X^{h,x}_n - \Delta X^{h,z}_n)^2 \mid \cF_n\right] \\
			%
			& = D_n^2 + 2 D_n \cdot h\left(b(X^{h,x}_n) - b(X^{h,z}_n)\right) + \E\left[(\Delta X^{h,x}_n - \Delta X^{h,z}_n)^2 \mid \cF_n\right]
		\end{align*}

		\noindent Lipschitzness of $b$ implies:
		%
		\[2 D_n \cdot h\left(b(X^{h,x}_n) - b(X^{h,z}_n)\right) \lesssim h D_n^2.\]
		%
		\noindent To handle the last term, first observe that if $(U^x, U^y)$ are coupled according to the optimal $W_2$ coupling of $\pi(\cdot \mid x)$ and $\pi(\cdot \mid y)$ then
		the Lipschitzness of $b, \pi$ imply:
		%
		\begin{align*}
			\E|b(x, U^x) - b(y, U^y)|^2 & \leq 2 \left[\E|b(x, U^x) - b(y, U^x)|^2 + \E |b(y, U^x) - b(y, U^y)|^2\right] \\
			%
			& \lesssim |x-y|^2, 
		\end{align*}
		%
		\noindent and similarly for $\sigma$. Hence the above suggests we couple $(X_t^{h,x}, X_t^{h,z})$ as follows. At every step, we have them share the same noise $\xi_n$ and
		we couple the controls $(U_n^x, U_n^y)$ according to the optimal $W_2$-coupling of $\pi(\cdot \mid X_n^{h,x}), \pi(\cdot \mid X_n^{h,z})$. Using $||Ax||_2 \leq ||A||_F ||x||_2$, this gives:
		%
		\begin{align*}
			& \E\left[(\Delta X^{h,x}_n - \Delta X^{h,z}_n)^2 \mid \cF_n\right]\\
			%
			& \leq h^2 \E \left|b(X_n^{h,x}, U_n^x) - b(X_n^{h,z}, U_n^z)\right|^2 + h \E \left|\sigma(X_n^{h,x}, U_n^x) - \sigma(X_n^{h,y}, U_n^y) \right|^2 \\
			%
			& \lesssim h|X_n^{h,x} -X_n^{h,y}|^2.\\ 
			%
			& = hD_n^2.
		\end{align*}
		%
		\noindent for $h \leq 1$. Combining with our above work, this implies:
		%
		\[\E \left[D_{n+1}^2 - D_n^2 \mid \cF_n\right] \lesssim h.\]
		%
		\noindent It follows by discrete Gr\"onwall's inequality that $\E D_n^2 \leq e^{C nh} D_0^2 \leq e^{CT}D_0^2$ and hence:
		% 
		\[\E| X^{h,x}_t -  X^{h, z_i}_t| \lesssim \epsilon.\]
		% 
		\noindent Combining our above results yields:
		%
		\[\sup_{t\in[0,T]} \sup_{x \in K}|P_t^h g_\pi(x) - P_t g_\pi(x)| \lesssim 2\epsilon + o(h). \]
		%
		\noindent where the $o(h)$ term does not depend on $\epsilon$,  giving the desired result.
	\end{proof}

	\noindent \textbf{Claim 2:} For every compact $K \subseteq \R^d$ and $T<\infty$:
			\[
		\sup_{x \in K}\int_0^T e^{-\rho t} |P_t g_\pi(x)|\,dt < \infty.
		\]

	\begin{proof}[Proof of Claim 2]
		We instead show the stronger statement:
		%
		\[\sup_{t \in [0, T]} \sup_{x \in K} |P_t g_\pi(x)| < \infty.\]
		%
		In Claim 1, we showed the map $x \mapsto P_tg_\pi(x)$ is Lipschitz in $x$, with Lipschitz constant
		uniform in $t$. Hence:
		%
		\[\sup_{x \in K} |P_t g_\pi(x)| \lesssim 1 + |P_t g_\pi(x_0)|,\]
		%
		\noindent for any fixed $x_0 \in K$. Thus it suffices to show $\sup_{t \in [0,T]} |P_t g_\pi(x_0)| < \infty$. As $g_\pi$ is Lipschitz
		we must have $|g_\pi(x)| \lesssim (1 + |x|)$ and hence:
		%
		\[\sup_{t \in [0,T]} |P_t g_\pi(x_0)| \lesssim 1 + \sup_{t\in [0,T]} \E_{x_0}|X^\pi_t| \leq 1 + \sqrt{\E_{x_0}\sup_{t\in [0,T]} |X^\pi_t|^2}, \]
		%
		\noindent which we show is finite in Claim 1 above.
	\end{proof}

	\noindent \textbf{Claim 3:} For every compact $K\subset \R^d$, if $\rho$ is sufficiently large

			\[
		\sup_{h\le 1}\sup_{x\in K}
		\int_T^\infty e^{-\rho t}|P_t^h g_\pi(x)|\,dt \to 0
		\quad \text{as } T\to\infty,
		\]
		and similarly
		\[
		\sup_{x\in K}
		\int_T^\infty e^{-\rho t}|P_t g_\pi(x)|\,dt \to 0
		\quad \text{as } T\to\infty.
		\]

	
	%
	
	\begin{proof}[Proof of Claim 3]
		Kolmogorov's backwards equations and linear growth give:
		%
		\[\frac{d}{dt}\E_x |X^\pi_t|^2 = \E_x \left[2 \langle b_\pi, X^\pi_t\rangle + |\sigma_\pi|^2 \right] \lesssim (1 + \E|X_t^\pi|^2).\]
		%
		for a constant $C$ not depending on $t$. Gr\"onwall's implies:
		%
		\[\E_x |X^\pi_t| \lesssim \left(1 + |x|\right) e^{Ct},\]
		%
		\noindent Hence for $\rho > C$, as \red{$g_\pi$ has linear growth} we see:
		%
		\[\sup_{x \in K}\int_T^\infty e^{-\rho t} |P_t g_\pi(x)| \lesssim \sup_{x \in K} (1 + |x|) \cdot \int_T^\infty e^{-(\rho - C) t} \; dt \to 0.\]
		%
		By considering the second moments of the discrete increments $X^h_n - X^h_{n-1}$, an identical argument handles the discrete case. 
	\end{proof}

	\noindent Using the above claims, we will prove the main result. Fix a compact set $K\subset \R^d$. Then for $x\in K$ and $T<\infty$, we have:
\[
V_\pi^h(x)-V_\pi(x)=I_1^h(x,T)+I_2^h(x,T)+I_3^h(x,T),
\]
where:
\begin{align*}
I_1^h(x,T)
&:=
\int_0^T e^{-\rho h\lfloor t/h\rfloor}\big(P_t^h g_\pi(x)-P_t g_\pi(x)\big)\,dt, \\
I_2^h(x,T)
&:=
\int_0^T \big(e^{-\rho h\lfloor t/h\rfloor}-e^{-\rho t}\big)P_t g_\pi(x)\,dt, \\
I_3^h(x,T)
&:=
\int_T^\infty e^{-\rho h\lfloor t/h\rfloor}P_t^h g_\pi(x)\,dt
-
\int_T^\infty e^{-\rho t}P_t g_\pi(x)\,dt.
\end{align*}

\noindent This separates $V_\pi^h - V_\pi$ into a semigroup error term $I_1^h$, a discounting error term $I_2^h$, and a tail error term $I_3^h$.
 We estimate each term uniformly on $K$. 

For the first term, by Claim 1,
\[
\sup_{x\in K}|I_1^h(x,T)|
\le
T \sup_{t\in[0,T]}\sup_{x\in K}|P_t^h g_\pi(x)-P_t g_\pi(x)|
\to 0.
\]

\noindent For the second term, since $0\le t-h\lfloor t/h\rfloor < h$, we have:
\[
|e^{-\rho h\lfloor t/h\rfloor}-e^{-\rho t}|
\le (e^{\rho h} - 1) e^{-\rho t}
\]
 \noindent Hence Claim 2 gives:
\[
\sup_{x\in K}|I_2^h(x,T)|
\le
(e^{\rho h} - 1) \cdot \sup_{x\in K} \int_0^T e^{-\rho t}|P_t g_\pi(x)|\,dt
\to 0.
\]

\noindent For the tail term, Claim 3 gives:
\[
\sup_{h\le1}\sup_{x\in K}|I_3^h(x,T)|\to 0
\qquad \text{as } T\to\infty.
\]

\noindent Now let $\varepsilon>0$. Choose $T$ so large that:
\[
\sup_{h\le1}\sup_{x\in K}|I_3^h(x,T)|<\varepsilon.
\]
Then choose $h$ sufficiently small so that
\[
\sup_{x\in K}|I_1^h(x,T)|+\sup_{x\in K}|I_2^h(x,T)|<\varepsilon.
\]
It follows that
\[
\sup_{x\in K}|V_\pi^h(x)-V_\pi(x)|<2\varepsilon
\]
for all sufficiently small $h$, which proves the result. 
\end{proof}

\noindent A standard localization argument gives the following:

\begin{corollary}
	Suppose $b, \sigma, g_\pi,$ and $\pi$ are only locally Lipschtiz. If we additionally 
	assume the global growth conditions:

	\begin{enumerate}[label=(\roman*)]
		\item $|b(x,u)| + |\sigma(x,u)| \leq C(1 + |u| + |x|).$
		\item $\int_U |u|^2 \; \pi(u \mid x) \leq C(1+ |x|^2)$.
		\item $|g_\pi(x)| \leq C(1 + |x|^2)$. 
	\end{enumerate}
	\noindent then the conclusion of Lemma \ref{lem:convergence_of_value_functions} still holds.
	\label{corr:conv_value_with_local_lipschitz}
\end{corollary}

\begin{proof}
	Fix a compact set $K$. For $R > \max_{x \in K} |x|$, let:
	%
	\[\tau_R := \inf \left\{t > 0 \; : \; |X_t^\pi| \geq R\right\}\]
	%
	\noindent be the first exit time of the ball of radius $R$. Analogously, define $\tau_R^h$ as the first
	exit time for the discrete chain. Our uniform $L^2$ bounds derived in Claim 1 above imply:
	%
	\[\P(\tau_R^h \leq T) = \P\left(\sup_{t \in [0,T]} |Y_t^h|^2 \geq R\right) \leq \frac{\E_x \sup_{t \in [0,T]} |Y_t^h|^2}{R^2} \lesssim\frac{e^{CT}}{R^2}. \]
	%
	\noindent uniformly in $h$. An analogous statement holds for $\tau_R$. \red{Moreover, the constants in these bounds are independent of the Lipschitz constants and $T$; they only depend on the growth conditions and $K$.} This suggests we can reasonably truncate
	our process without too much error, and apply Lemma \ref{lem:convergence_of_value_functions} to the truncated processes. 
	
	More formally, let $\phi_R: \R^d \to \overline{B}(0, R)$ be the orthogonal projection onto the closed ball of radius $R$. Define new coefficients and policy:
	%
	\[b_R(x, u) := b(\phi_R(x), u), \quad \sigma_R(x, u) := \sigma(\phi_R(x), u), \quad \pi_R(\cdot \mid x) := \pi( \cdot \mid \phi_R(x)).\]
	%
	\noindent For fixed $R$, these are globally Lipschitz, and satisfy all the conditions of Lemma \ref{lem:convergence_of_value_functions}. Let $V_\pi^{h, R}$ and $V_\pi^R$ denote the
	value functions corresponding to these new coefficients/policy for the discrete and continuous processes respectively. It suffices then to show that the truncation error tends to zero uniformly in $h, x$ as $R \to \infty$:
	%
	\[\sup_{h \leq 1} \sup_{x \in K} |V_\pi^h(x) - V_\pi^{h,R}(x)| \to 0, \quad  \sup_{x \in K} |V_\pi(x) - V_\pi^{R}(x)| \to 0\]
	%
	\noindent Note by the growth bound on $g_\pi$ and Cauchy-Schwarz:
	%
	\begin{align*}
		|V_\pi(x) - V_\pi^R|^2 & \lesssim \left(\int_0^\infty e^{-\rho t} \E_x \left[(1 + |X_t^\pi|^2) \cdot 1\{\tau_R \leq t\} \right]\; dt \right)^2 \\
		%
		& \lesssim \int_0^\infty e^{-2\rho t} (1 + \E_x|X_t^\pi|^4) \;dt \cdot\int_0^\infty \P(\tau_R \leq t) \; dt\\
		%
		& \red{\lesssim \frac{(1 + |x|^2)(1 + |x|^4)}{R^2} \int_0^\infty e^{-(\rho - C)t} dt} 
	\end{align*}
	\red{todo; we can bound the first using Kolmogorovs backwards equations as before, and the second we bound using our above statement. Analogously, we can bound the discrete chain.}
\end{proof}


%%%%%%%%%%%%%%%%%%%%%%%%%%%%%%
% Convergence of Optimal Policies
%%%%%%%%%%%%%%%%%%%%%%%%%%%%%%
\subsection{Convergence of Optimal Policies}

\noindent Now, we handle the convergence of the relevant optimal policies.

 For the discrete chain $X^h_n$, if  the optimal value function \red{$V^h \in C^2$}, it satisfies:
%
\begin{align*}
	Q^h_\gamma(x,u) & = h \ell(x,u) + e^{-\rho h} \E_\xi V^h_\pi(x + h b(x,u) + \sqrt{h} \sigma(x,u) \xi), \\
					& = V^h + h \left[\ell + V_x^h \cdot b + \frac{\sigma^2}{2} \cdot V_{xx}^h - \rho V^h\right] + o(h).
\end{align*}
%
\noindent Moreover, Bellman's equation \eqref{eqn:entropy_regularized_bellman_eqn} then implies:
%
\[\rho V^h(x) = \inf_{\mu \in \cP(U)} \left\{\int_U \left[\ell + V_x^h \cdot b + \frac{\sigma^2}{2} \cdot V_{xx}^h + \tau \log\left(\frac{d\mu}{d\nu}\right) + o(1)\right] \mu(du)\right\}. \]
%

\noindent Since we have uniform convergence of these to the limiting value function, some simple analysis 
of the discrete Bellman equation \eqref{eqn:entropy_regularized_bellman_eqn} suggests the form of the 
limiting Bellman equation. Namely:
%
\[\rho V(x) = \inf_{\mu \in \cP(U)} \left\{\int_U \left[\ell + V_x\cdot b + \frac{\sigma^2}{2} \cdot V_{xx} + \tau \log\left(\frac{d\mu}{d\nu}\right) \right] \mu(du)\right\}. \]
%
\noindent which recovers equation 14 in \cite{WangZariphopoulouZhou2020}.

\begin{lemma}[Convergence of Optimal Policy]
	\red{Under some conditions} we have \red{pointwise convergence of the optimal policies}:
	%
	\[\pi^h(u \mid x) \; d\nu = Z^h(x)^{-1}\exp\left(-\frac{Q^h_\gamma(x,u)}{h\tau} \right) \to Z(x)^{-1}\exp\left(-\frac{Q_\gamma(x,u)}{\tau} \right) \; d\nu = \pi^*(u \mid x)\]
	%
\end{lemma}

\begin{proof}

\end{proof}










\newpage 
%%%%%%%%%%%%%%%%%%%%%%%%%%%%%%%%%%%%%%%%%%%%%%%%%%%%
% LQ Optimal Control 
%%%%%%%%%%%%%%%%%%%%%%%%%%%%%%%%%%%%%%%%%%%%%%%%%%%%
\subsection{Linear-Quadratic Optimal Control}

\red{We should show the conditions/assumptions we make above naturally hold under the simple case of the LQR. 
	Then we should just briefly review the results of \cite{WangZariphopoulouZhou2020} in terms of what the limit looks like, and what our corresponding 
	discrete approximations look like. }\\




\zack{Can we actually compute what the optimal policy is in the discrete chain case?}

%%%%%%%%%%%%%%%%%%%%%%%%%%%%%%%%%%%%%%%%%%%%%%%%
% References
%%%%%%%%%%%%%%%%%%%%%%%%%%%%%%%%%%%%%%%%%%%%%%%%
\newpage
% include hyperlinks for references?
\bibliographystyle{plainnat}
\bibliography{references}




%%%%%%%%%%%%%%%%%%%%%%%%%%%%%%%%%%%%%%%%%%%%%%%%
% Appendix
%%%%%%%%%%%%%%%%%%%%%%%%%%%%%%%%%%%%%%%%%%%%%%%%
\newpage 
\section{Appendix}

We start by proving the optimal policy for the entropy-regularized 
control problem takes the form \eqref{eqn:optimal_policy_entropy_regularized}. As before, let:
%
\[Z(x) := \int_U \exp\left(-\frac{Q_\gamma(x,u)}{\tau}\right) \; d\nu(u)\]
%
\noindent denote the partition function.
%
\begin{lemma}[Optimal Policy for Entropy-Regularized Control]
	If $0 < Z(x) < \infty$ then: 
		\[\pi^*(u \mid x) \; d\nu = Z(x)^{-1}\exp\left(-\frac{Q_\gamma(x,u)}{\tau} \right), \quad V(x) = -\tau \log(Z(x)).\]
	%	
	is the optimal policy for the entropy-regularized Bellman equation \eqref{eqn:entropy_regularized_bellman_eqn}:
	%
	\[L(\mu) := \int_U \left[Q_\gamma(x,u) + \tau \log\left(\frac{d\mu}{d\nu}\right)\right] \;  d\mu.\]
	%
\end{lemma}
%
\begin{proof}
	Suppose $\mu \in \mathcal{P}(U)$ is any other policy. We can assume $\mu \ll \nu$ as otherwise 
	$KL(\mu \mid \nu) = \infty$, so clearly $\mu$ cannot be optimal. Note:
	%
	\[Q_\gamma(x,u) = -\tau \log (Z(x)) - \tau \log\left(\frac{d\pi^*}{d\nu}\right), \quad L(\pi^*(\cdot \mid x)) = -\tau \log(Z(x)). \]
	%
	\noindent Hence: 
	%
	\begin{align*}
		L(\mu) - L(\pi^*(\cdot \mid x)) & = \int_U \left[Q_\gamma(x,u) + \tau \log\left(\frac{d\mu}{d\nu}\right)\right] \; \mu(du) + \tau \log(Z(x))\\
		& = \tau \int_U \left[\log\left(\frac{d\mu}{d\nu}\right) - \log\left(\frac{d\pi^*}{d\nu}\right)\right] \; \mu(du)\\
		& = \tau KL(\mu \mid \pi^*(\cdot \mid x))\\
		& \geq 0
	\end{align*}
	%
	\noindent where the above is infinite if $\mu$ is not absolutely continuous with respect to $\pi^*(\cdot \mid x)$. Hence $\pi^*(\cdot \mid x)$ is optimal.
\end{proof}

\end{document}
